\documentclass[11pt]{article}

\usepackage[T1]{fontenc}
\usepackage[utf8]{inputenc}
\usepackage{lmodern}
\usepackage[margin=1in]{geometry}
\usepackage{amsmath,amssymb,amsthm,mathtools,bm}
\usepackage{booktabs}
\usepackage{enumitem}
\usepackage[dvipsnames]{xcolor}
\usepackage{microtype}
\usepackage[round,authoryear]{natbib}
\usepackage[colorlinks=true,linkcolor=MidnightBlue,citecolor=MidnightBlue,urlcolor=MidnightBlue]{hyperref}

\setlist{itemsep=2pt,topsep=4pt}

\theoremstyle{plain}
\newtheorem{proposition}{Proposition}
\newtheorem{target}{Target result}
\renewcommand{\thetarget}{\Alph{target}}

\newcommand{\E}{\mathbb{E}}
\newcommand{\R}{\mathbb{R}}
\newcommand{\cZ}{\mathcal{Z}}
\newcommand{\cI}{\mathcal{I}}
\newcommand{\QJ}{Q_J}
\newcommand{\norm}[1]{\lVert #1 \rVert}
\newcommand{\ind}{\mathbf{1}}
\DeclareMathOperator*{\argmin}{arg\,min}
\newcommand{\dfsp}{\textnormal{D-FSP}}
\newcommand{\sfsp}{\textnormal{S-FSP}}
\newcommand{\fsppo}{\textnormal{FSP-PO}}

\title{Distributional Few-Shot Personalization\\for Human-Anchored LLM Post-Training\\[6pt]
\large Research proposal}
\author{LLM-humanchoice project}
\date{September 2026}

\begin{document}
\maketitle

\begin{abstract}
LLM judges are cheap and can report a full distribution over preference outcomes, but they do not
agree with human raters, and human labels are expensive. We propose to treat the judge as a
distribution-valued black-box predictor and to correct it toward human preferences using only a
few human labels. The approach builds on the few-shot personalization framework of
\citet{li2026personalizing}. The project has three parts:
\begin{itemize}
  \item \emph{Distributional few-shot personalization} (\dfsp): the judge's win probability serves as
        an offset, and a kernel estimate of the judge-to-human discrepancy, localized on summaries
        of the judge distribution, corrects it.
  \item \emph{A label-acquisition rule} that decides which comparisons humans should label.
  \item \emph{Post-training (\fsppo)}: the corrected human preference probabilities serve as soft
        targets for DPO-style training of an open-weight model.
\end{itemize}
We aim to show three things: when the judge distribution carries information beyond its win
probability, how few human labels are needed, and how the estimation error carries over to the
post-trained policy. All experiments reuse an existing human-labeled dataset, revealing only a
small fraction of its labels.
\end{abstract}

\section{Introduction}
\label{sec:intro}

LLMs are widely used as judges for preference annotation and post-training
\citep{bai2022constitutional,zheng2023judging,lee2024rlaif}. For a prompt $x$ and two responses
$y_A,y_B$, a judge queried by logit elicitation or repeated sampling returns a distribution
$\QJ(\cdot\mid X)$ over preference outcomes, $X=(x,y_A,y_B)$. This distribution is informative but
systematically different from the human preference we want to align to. Human labels are far more
expensive, so in practice we have judge outputs on many comparisons and human labels on only a
few. We ask how to combine the two into a human-anchored preference target, and then how to train a
policy against that target.

\paragraph{Starting point.}
\citet{li2026personalizing} study few-shot personalization of a black-box predictor in
nonparametric regression. Their procedure has three steps:
\begin{enumerate}
  \item apply a local smoothing to the pretrained predictor;
  \item estimate the discrepancy between the target and the smoothed predictor with a kernel
        smoother on a few target labels;
  \item choose the amount of smoothing on validation data.
\end{enumerate}
The candidate set in step~3 includes the target-only estimator, so an uninformative predictor is
automatically ignored. They also give a variance-driven rule for where to collect labels, and prove
a minimax rate that improves on target-only estimation when the discrepancy is simpler than the
target. Their setting is scalar regression, which includes binary classification. It does not
cover graded black-box outputs or preference post-training.

\paragraph{This proposal.}
We extend this idea in three ways.
\begin{enumerate}
  \item \textbf{Distributional few-shot personalization (\dfsp).} The judge's win probability is the
        offset. The discrepancy is localized on richer summaries of a graded judge distribution
        (tie mass, order inconsistency). This lets the correction use information that a scalar
        probability discards.
  \item \textbf{Label acquisition.} A design rule for which comparisons to send to humans, derived
        from the risk of the estimator. It favors regions where humans disagree and where the
        judge-to-human discrepancy changes quickly.
  \item \textbf{Post-training (\fsppo).} The corrected probabilities are soft targets for DPO-style
        training. We show how the estimation error bounds the excess preference risk of the
        trained policy.
\end{enumerate}

\paragraph{Related work.}
Soft AI preference labels are already used for reward modeling and preference optimization
\citep{lee2024rlaif,furuta2024gapo}, as are strength-aware and distributional preference models
\citep{imai2026mixdpo,li2024dprm}. Other work combines few human labels with AI feedback through
targeted correction \citep{xu2025rlthf}, a learned linear map from judge outputs to human judgments
\citep{vandenburg2025aligning}, or a jointly weighted Bradley--Terry likelihood whose weight can fall
back to human-only estimation \citep{cai2026dial}. Active selection of pairs for human labeling has
also been studied \citep{shen2025active}. Our contribution is none of these components on its own.
It is a nonparametric, adaptive correction of a \emph{distribution-valued} judge, with guarantees
that carry through to post-training.

\section{Setup}
\label{sec:setup}

\paragraph{Data.}
We use an existing dataset of pairwise comparisons $X=(x,y_A,y_B)$ with binary human labels
$Y\in\{0,1\}$, where $Y=1$ means $A$ is preferred. The judge is run on all comparisons. During
training only $n$ human labels are revealed. A held-out test set, split by prompt, is used for
evaluation. No new human annotation is collected.

\paragraph{Judge protocol.}
A binary verdict carries no information beyond its win probability $p_J$, so we elicit a graded
distribution instead. The judge chooses among five labels, $\{A\gg B,\ A>B,\ \text{tie},\ B>A,\
B\gg A\}$. We read the label probabilities from its output logits, once with $A$ shown first
($\QJ^{AB}$) and once with the order reversed ($\QJ^{BA}$). Let
$w(Q)=Q(A\gg B)+Q(A>B)+\tfrac12Q(\text{tie})$. We summarize each comparison by three quantities:
\[
  p_J=\tfrac12\{w(\QJ^{AB})+w(\QJ^{BA})\},\qquad
  t_J=\tfrac12\{\QJ^{AB}(\text{tie})+\QJ^{BA}(\text{tie})\},\qquad
  o_J=|w(\QJ^{AB})-w(\QJ^{BA})| .
\]
These are the win probability, the tie mass and the order inconsistency. The last two are not
functions of $p_J$.

\paragraph{Target.}
Let $Z=(p_J,t_J,o_J)\in\cZ$, a fixed polytope in $[0,1]^3$. When content categories are available,
$Z$ is used within each category. The target is the human preference probability
\[
  p_H(z)=P(Y=1\mid Z=z),
\]
and an estimator $\hat p$ is evaluated by
\[
  R(\hat p)=\E\{\hat p(Z)-p_H(Z)\}^2 .
\]
On held-out comparisons with one label each, the Brier score equals $R(\hat p)$ plus a constant
that does not depend on the method. Methods can therefore be compared directly.

\section{Methods}
\label{sec:methods}

\subsection{Method I: Distributional few-shot personalization (\dfsp)}
\label{sec:dfsp}

\paragraph{Decomposition.}
We write the human target as the judge's prediction plus a discrepancy,
\[
  p_H(z)=p_J(z)+\delta(z).
\]
Here $p_J(z)$ is the first coordinate of $z$. If the judge is informative, $\delta$ is simpler than
$p_H$: it is smaller, or smoother, or both. It can then be estimated from far fewer labels than
$p_H$ itself.

\paragraph{Local smoothing.}
Following \citet{li2026personalizing}, for $\theta=(\theta_1,\theta_2)$ with $\theta_1\ge0$ and
$0<\theta_2\le1$, and a center $z$, define the truncated offset
\[
  \omega_{\theta,z}(z')=p_J(z)+\min\bigl\{|p_J(z')-p_J(z)|,\ \theta_1\norm{z'-z}^{\theta_2}\bigr\}\,
  \operatorname{sgn}\{p_J(z')-p_J(z)\}.
\]
This operation limits how much the offset may vary around $z$, and $\theta$ acts as a borrowing
dial. When $\theta_1=0$, $\omega_{\theta,z}\equiv p_J(z)$, and the estimator below reduces to a purely
human kernel estimate. When $\theta_1$ is large, the full offset is used.

\paragraph{Estimator.}
Let $\cI_n$ index the revealed labels and let $K_h(z',z)=\ind\{\norm{z'-z}_\infty\le h\}$. Define
\[
  \hat\delta_{\theta,h}(z)=\frac{\sum_{i\in\cI_n}K_h(Z_i,z)\{Y_i-\omega_{\theta,z}(Z_i)\}}
                                 {1\vee\sum_{i\in\cI_n}K_h(Z_i,z)},
  \qquad
  \hat p_{\theta,h}(z)=\Pi_{[0,1]}\{p_J(z)+\hat\delta_{\theta,h}(z)\},
\]
where $\Pi_{[0,1]}$ clips to $[0,1]$. The pair $(\theta,h)$ is chosen from a grid by validation
Brier score on a held-out part of the revealed labels. Because the grid contains $\theta_1=0$, a
judge that does not help is automatically switched off.

\paragraph{Scalar versus distributional localization.}
Scalar FSP (\sfsp) localizes on $p_J$ alone. \dfsp{} localizes on $Z=(p_J,t_J,o_J)$. The two
procedures target different quantities, $\E(Y\mid p_J)$ and $\E(Y\mid Z)$. The comparison between
them decomposes exactly as
\[
  R_X(\hat p_{\mathrm S})-R_X(\hat p_{\mathrm D})
  =G-\{R(\hat p_{\mathrm D})-R(\hat p_{\mathrm S})\},
  \qquad
  G=\E\{\E(Y\mid Z)-\E(Y\mid p_J)\}^2 .
\]
Here $R_X$ is the error against the per-comparison probability $P(Y=1\mid X)$. Each $R(\cdot)$ is
measured against its own target, so the braced term is the extra estimation cost of localizing in
three dimensions instead of one. The judge distribution helps exactly when its information gain
$G$ exceeds that extra cost. $G$ is a variable-importance quantity
\citep{williamson2021nonparametric,williamson2023general}, and we test it directly
(Section~\ref{sec:experiments}).

\subsection{Method II: Choosing which comparisons to label}
\label{sec:acquisition}

We choose the $n$ comparisons to label from the judged pool. Let $q$ be the density of $Z$ in the
pool, $\sigma^2(z)=p_H(z)\{1-p_H(z)\}$ the human label noise, and $\gamma_1(z)$ the local
H\"older constant of the discrepancy $\delta$. Suppose labels are drawn with density $p$ and the
kernel uses a locally chosen bandwidth. Minimizing the resulting bias--variance upper bound on the
risk over $p$ gives the design
\[
  p^{*}(z)\;\propto\;q(z)^{\frac{2\gamma_2+d}{4\gamma_2+d}}\;
  \gamma_1(z)^{\frac{2d}{4\gamma_2+d}}\;\sigma(z)^{\frac{4\gamma_2}{4\gamma_2+d}} .
\]
In words, label comparisons that are \emph{common}, on which \emph{humans are split}, and where
the judge-to-human discrepancy \emph{varies}. A constant bias of the judge is cheap to learn
however large it is. Labels should not simply go where the judge is unsure.

In practice we proceed as follows:
\begin{enumerate}
  \item Reveal a small uniform pilot sample.
  \item Fit \dfsp{} on the pilot sample.
  \item Estimate $\sigma$ from the pilot fit.
  \item Estimate $\gamma_1$ from local differences of the fitted discrepancy, at a scale larger than
        the pilot bandwidth.
  \item Plug the estimates into $p^{*}$, mix the result with $q$ for robustness, and sample the
        remaining labels from the pool accordingly.
\end{enumerate}
For H\"older classes no design can improve the minimax \emph{rate} \citep{castro2005faster}. The
goal of this rule is a better constant, that is, fewer labels for the same accuracy.

\subsection{Method III: Post-training against the personalized target (\fsppo)}
\label{sec:post-training}

Let $\pi_{\mathrm{ref}}$ be the initial model and $\pi_\vartheta$ the model being trained. Following
DPO \citep{rafailov2023dpo}, define
\[
  d_\vartheta(X)=\beta\Bigl[\log\frac{\pi_\vartheta(y_A\mid x)}{\pi_{\mathrm{ref}}(y_A\mid x)}
  -\log\frac{\pi_\vartheta(y_B\mid x)}{\pi_{\mathrm{ref}}(y_B\mid x)}\Bigr].
\]
We train on all judged pairs with the soft-label objective
\[
  \widehat R(\vartheta;\hat p)=-\frac1N\sum_{i=1}^{N}
  \Bigl[\hat p_i\log\sigma\{d_\vartheta(X_i)\}+(1-\hat p_i)\log\sigma\{-d_\vartheta(X_i)\}\Bigr],
\]
where $\hat p_i=\hat p(Z_i)$ comes from \dfsp. The loss itself is standard
\citep{lee2024rlaif,furuta2024gapo}. What is new is the target $\hat p$ and the guarantee it carries.

Let $R(\vartheta;p)$ denote the population version of this objective, and let $p_X=P(Y=1\mid X)$.

\begin{proposition}[From estimation error to post-training error]
\label{prop:bridge}
Suppose $|d_\vartheta|\le D$, and $\hat\vartheta$ minimizes $R(\cdot;\hat p)$ up to an error
$\varepsilon$ (optimization plus generalization). Then
\[
  R(\hat\vartheta;p_X)-\inf_\vartheta R(\vartheta;p_X)\;\le\;\varepsilon+2D\,R_X(\hat p)^{1/2}.
\]
Suppose in addition the class $\{d_\vartheta\}$ is convex. Then the bound improves to
$2\varepsilon+2R_X(\hat p)/c_D$, where $c_D=\sigma(D)\{1-\sigma(D)\}$.
\end{proposition}

\noindent\emph{Proof idea.} The objective is linear in the target:
$\ell(d;p)-\ell(d;p')=-(p-p')\,d$. Replacing $p_X$ by $\hat p$ therefore shifts every candidate's
risk by at most $D\,\E|\hat p-p_X|$, which gives the first bound. For a convex class, the logistic
loss grows quadratically around its minimizer, which gives the second.

The proposition turns the human-label cost of estimating $p_H$ into a bound on the policy's excess
preference risk. $R_X$ includes the part of human preference not explained by $Z$, which no
amount of human labels removes.

\section{Theoretical goals}
\label{sec:theory}

The theory adapts \citet{li2026personalizing} to bounded preference probabilities and to a
localization variable built from the judge distribution.

\paragraph{Assumptions.}
\begin{enumerate}[label=(\roman*)]
  \item $p_H$ is H\"older with parameters $(\theta_1^{*},\theta_2^{*})$.
  \item Preferences are nondegenerate: $p_H\in[\kappa,1-\kappa]$.
  \item Labeled points have a density bounded above and below on a regular subset of $\cZ$.
  \item The bandwidth satisfies $nh^d\gtrsim\log n$.
\end{enumerate}
We say the discrepancy has complexity $\gamma=(\gamma_1,\gamma_2)$ if, after smoothing, it is
H\"older with constant $\gamma_1$ and exponent $\gamma_2$.

\begin{target}[Upper bound]
For fixed $\theta$ and a suitable bandwidth $h$,
\[
  R(\hat p_{\theta,h})\;\lesssim\;\gamma_1^{\frac{2d}{2\gamma_2+d}}\,
  n^{-\frac{2\gamma_2}{2\gamma_2+d}}+\frac1n .
\]
\end{target}

\begin{target}[Lower bound]
Take a fixed judge that lies strictly inside the probability range and strictly inside the
smoothness class of $p_H$, and suppose $\theta_2^{*}\le\gamma_2\le1$. Then the rate of Target~A
cannot be improved by any estimator and any labeling design, including sequential designs.
\end{target}

\begin{target}[Adaptation and no harm]
With $(\hat\theta,\hat h)$ chosen by validation,
\[
  R(\hat p_{\hat\theta,\hat h})\le C\min_{\theta,h}R(\hat p_{\theta,h})+O\{(\log n)^2/n\}.
\]
In particular, \dfsp{} is never much worse than the human-only estimator.
\end{target}

\paragraph{Consequences.}
A good judge has a discrepancy with small $\gamma_1$ or large $\gamma_2$, so it lowers the number of
human labels needed. Combining Targets~A--B with Proposition~\ref{prop:bridge} gives a post-training
statement: under realizability, the policy's excess preference risk decays at the same rate. That
rate is governed by the complexity of the judge-to-human discrepancy, not by the complexity of
human preference itself. For the design of Section~\ref{sec:acquisition}, we aim to show that the
plug-in rule attains the risk bound of the oracle design up to the cost of the pilot sample.

\section{Experiments}
\label{sec:experiments}

All experiments reuse the human-labeled comparison dataset provided for this project. We run the
judge on every comparison and use $n$ human labels, with $n/N\in\{1,2,5,10,20\}\%$. A test set split
by prompt is used only for evaluation.

\paragraph{E1: Estimation.}
We compare \dfsp{} with the following estimators:
\begin{itemize}
  \item the judge's win probability;
  \item a kernel estimator using human labels only;
  \item Platt scaling and isotonic regression of $p_J$;
  \item the linear map of \citet{vandenburg2025aligning};
  \item a GAM and boosted trees on $Z$;
  \item \sfsp.
\end{itemize}
We report Brier score, log-loss and calibration error as functions of $n$. We also run \dfsp{} with
only the tie probability or only the order dependence added to $p_J$.

\paragraph{E2: Does the judge's distribution add information?}
We estimate $G$ from Proposition~\ref{prop:value} and test $G=0$ with the method of
\citet{williamson2023general}. The test is first checked on simulated data with known $G$.

\paragraph{E3: Label selection.}
We compare the rule of Section~\ref{sec:selection} with the following alternatives:
\begin{itemize}
  \item random selection;
  \item selection by judge uncertainty;
  \item Fisher-information selection \citep{shen2025active};
  \item selection of likely judge errors, as in \citet{xu2025rlthf}.
\end{itemize}
We run this comparison on simulated data with known $p_H$ and $\delta$, and on the real data.

\paragraph{E4: Unreliable judge.}
We corrupt the judge's probabilities in four ways: adding a bias, making them overconfident,
reversing them in part of $\cZ$, and adding noise. We check that the selected $\theta_1$ decreases
and that \dfsp{} approaches the human-only estimator.

\paragraph{E5: Post-training.}
We fine-tune a small open-weight model with LoRA \citep{hu2022lora} under the following training
targets:
\begin{itemize}
  \item the $n$ human labels only;
  \item hard judge labels;
  \item judge probabilities, trained with soft DPO and with GAPO \citep{furuta2024gapo};
  \item linear-map estimates;
  \item \dfsp{} estimates;
  \item all human labels, as a reference.
\end{itemize}
We report the held-out human log-loss of $\sigma(d_\vartheta)$, which is the quantity in
Proposition~\ref{prop:bridge}. We also report the win rate against $\pi_{\mathrm{ref}}$, scored by a
reward model trained on all human labels \citep{gao2023scaling} and by a second judge. Finally, we
report the KL divergence to $\pi_{\mathrm{ref}}$ and the response length.

\paragraph{Possible outcomes.}
If $G$ is close to zero, \dfsp{} should perform like \sfsp, and the value of the method lies in the
correction and the label selection alone. If the judge is poor, \dfsp{} should perform like the
human-only estimator. Both outcomes will be reported.

\section{Expected contributions}
\label{sec:contributions}

\begin{enumerate}
  \item \textbf{Method.} \dfsp{} corrects a distribution-valued LLM judge toward human preferences
        with few labels. It is adaptive: it uses the judge when the judge helps and falls back to
        human-only estimation when it does not.
  \item \textbf{Label design.} A rule for which comparisons to send to humans, targeting regions
        where humans disagree and the judge's error varies.
  \item \textbf{Post-training.} A pipeline from few human labels to a post-trained policy, with a
        bound linking estimation error to the policy's excess preference risk, and experiments
        measuring policy quality against the human-label budget.
\end{enumerate}

\noindent The main hypothesis can also fail in informative ways. If the graded judge distribution
adds nothing beyond $p_J$ ($G=0$), \dfsp{} should reduce to scalar FSP. If the judge is misleading,
\dfsp{} should revert to human-only estimation. Both outcomes will be reported.


\bibliographystyle{plainnat}
\bibliography{references}

\end{document}
